When we write a program in our favorite programming language, we fill a file with a collection of symbols that we've typed out on a keyboard.
This is one of the beauties of programming: looking past the bells and whistles provided by editors and programming tools, a program is just a stream of characters.
At some point in our programming workflow, we need to verify that what we've typed up works as expected, so we \textit{run} our program.
In some editors, there is literally \enquote{play} button, but in many other cases this means opening a terminal and typing out a few commands.
In either case, we're running a \textit{different} program in order to run the program we've written.
Our goal is to understand what's going on here:
\textit{
  What is this program doing?
  How does it do it?
  What sorts of data structures does this program use?
}
Our intuition should tell us this program is probably doing something fairly complicated.\footnote{Just think of how hard it is to follow directions, and how often we wish we read \textit{all} the instructions before starting some task.}
So, our basic question is: \textit{How do we get from a stream of symbols to the output of the program?}

\begin{figure}
  \hspace{-1.5cm} % ugh
  \begin{tikzpicture}[
      interpbox/.style={
        rectangle,
        rounded corners,
        draw,
        thick,
        minimum width=6cm,
        minimum height=1cm,
      }
    ]
    \draw[thick, rounded corners] (2,1) rectangle (9, -12.25);
    \node (chars) {\textit{stream of characters}};
    \node[interpbox] (lex) [right=of chars] {Lexical Analysis};
    \node (tokens) [below=of lex] {\textit{stream of tokens}};
    \node[interpbox] (par) [below=of tokens] {Syntactic Analysis (Parsing)};
    \node (ast) [below=of par] {\textit{abstract syntax tree (AST)}};
    \node[interpbox] (sem) [below=of ast] {Semantic Analysis (Type Checking)};
    \node (ir) [below=of sem] {\textit{well-typed AST/intermediate representation}};
    \node[interpbox] (eval) [below=of ir] {Evaluation};
    \node (input) [left=of eval] {\textit{input}};
    \node (output) [right=of eval] {\textit{output}};
    \draw[->, thick] (chars.east) -- (lex.west);
    \draw[thick] (lex.south) -- (tokens.north);
    \draw[->, thick] (tokens.south) -- (par.north);
    \draw[thick] (par.south) -- (ast.north);
    \draw[->, thick] (ast.south) -- (sem.north);
    \draw[thick] (sem.south) -- (ir.north);
    \draw[->, thick] (ir.south) -- (eval.north);
    \draw[->, thick] (input.east) -- (eval.west);
    \draw[->, thick] (eval.east) -- (output.west);
    \node at (7.7, -12) {\texttt{INTERPRETER}};
  \end{tikzpicture}
  \centering
  \caption{
    Visualization of the interpretation pipeline
  }
  \label{fig:nested-progs}
\end{figure}

\todo{expections of the chapters}

\section{The Interpretation Pipeline}

\todo{outline of the text}

\todo{fill this out}
